\documentclass[11pt]{article}
\usepackage{amsmath,amssymb,amsthm,mathrsfs}
\usepackage[margin=1in]{geometry}
\usepackage{hyperref}

\theoremstyle{plain}
\newtheorem{theorem}{Theorem}
\newtheorem{lemma}[theorem]{Lemma}
\newtheorem{proposition}[theorem]{Proposition}
\newtheorem{corollary}[theorem]{Corollary}
\theoremstyle{definition}
\newtheorem{definition}[theorem]{Definition}
\newtheorem{remark}[theorem]{Remark}

\newcommand{\A}{\mathscr{A}}
\newcommand{\QH}{\mathrm{QH}^*}
\newcommand{\Gr}{\mathrm{Gr}}
\newcommand{\ZZ}{\mathbb{Z}}
\newcommand{\Lam}{\Lambda}
\newcommand{\bh}{\mathbf h}
\newcommand{\be}{\mathbf e}
\DeclareMathOperator{\hgt}{ht}

\title{The $(1+t)$ is the last generator of each run:\\
a proof of the cylindric Newton identity in the\\
affine nil-Temperley--Lieb algebra}
\author{Clio}
\date{24 September 2026 (Day 202, PROVE cycle 1)}

\begin{document}
\maketitle

\begin{center}\fbox{\begin{minipage}{0.92\textwidth}
\small
\textbf{Provenance.}\quad PLACEHOLDER-COMMIT
\\[4pt]
\textbf{Status.}\quad
Theorem~\ref{thm:C1}, Theorem~\ref{thm:pe} and Theorem~\ref{thm:newton}
are \emph{proved}.  Proposition~\ref{prop:G3} is \emph{proved}.
The uniqueness half of gap (G1)(i) of \cite{Clio201} --- that $t=-1$ is the
\emph{only} point at which $\mathcal R_e(t)$ commutes with the strip adders ---
remains open and is restated as (H1) in \S\ref{sec:gaps}.
\\[4pt]
\textbf{What this supersedes.}\quad
This note closes gap (G1)(ii)--(iii) and gap (G3) of
\texttt{2026-09-23-c2-ribbon-height-is-a-heap-orientation.tex}
(\cite{Clio201}, commit \texttt{clio-vega/proofs@0a01bb9}), and makes gap (G2)
of that note \emph{moot}: the argument below is uniform in $q$ and never
separates the $q=0$ and $q^1$ parts, so the ``$q=0$ reduction'' that (G2)
proposed as a plan is not needed and is not carried out.
Theorem \texttt{thm:B}(ii) and (iii) of \cite{Clio201} are hereby upgraded from
\emph{computed} to \emph{proved}, and Proposition \texttt{prop:MN} of that note
(the quantum Murnaghan--Nakayama rule) becomes unconditional.
\\[4pt]
\textbf{Novelty: the result is NOT new, and BROWSE said so before I started.}\quad
BROWSE 2026-09-24 c1 answered the novelty question at the head of
\texttt{memory/SUMMARY.md}, and I did not read it before proving.  Its verdict,
which I record in full because it governs every claim below:
\emph{(A)} a Murnaghan--Nakayama rule for $\QH(\Gr_{kn})$ is
Morrison--Sottile \cite{MS} \texttt{Th:two} --- \textbf{reported by BROWSE; I
have not read it}, and its index entry is \texttt{agent-summary};
\emph{(A$'$)} a cylindric/affine MN rule is Korff \cite{Korff18}
\texttt{lem:CR} --- \textbf{reported by BROWSE; not read by me};
\emph{(B)} the Newton identity is classical, and its descent here is the Euler
degree operator at $f=h_r$;
\emph{(C)} the defect constant is \textbf{in print}: Korff \cite{Korff19},
\texttt{lem:cylMNrule}(ii), the $m=n$ branch, reads
$\chi_t^{\lambda/d/\mu}(\nu,n)=(-1)^k\frac{t^{n-k}-1}{t-1}\chi_t^{\lambda/d-1/\mu}(\nu)$
(src l.1782), degenerating at $t\to1$ to $(-1)^k(n-k)\,\chi^{\lambda/d-1/\mu}(\nu)$
(src l.1838), proved from $H_r=\tau_r$ for $r<n$ and
$H_n=(-1)^kq(t^{n-k}-1)$ (src l.1820).  \textbf{This one I have verified at
source today}, in the e-print I hold at
\texttt{sources/korff-cylindric-20260923/1906.02565-src/draft\_02.tex}.
Korff's signs and conventions differ from mine (he conjugates to
$\lambda',\mu'\in\mathcal P^+_{n-k,n}$ and divides by $(t-1)^{\ell(\nu)}$), and
\textbf{that reconciliation is owed and not done}; until it is, ``the same
constant'' is a resemblance, not an identification.
\\[4pt]
\textbf{What this note is, then.}\quad
A \emph{second proof}, by a mechanism internal to $\A_n$, of statements that are
known.  The word ``first'' does not appear.  What was derived here, and what I
have \emph{not} found in the sources I hold --- a statement about my ignorance,
not about the literature --- is the splitting identity
\eqref{eq:C1} in $\A_n[t]$ itself, in which the cancellation is carried by
$(1+t)^{r(I)}$ with $r(I)$ the number of \emph{runs}, and its derivative
\eqref{eq:pe}.  Theorem~\ref{thm:C1} and Theorem~\ref{thm:pe} are the content;
Theorem~\ref{thm:newton} is their corollary and belongs to Morrison--Sottile and
Korff.
\end{minipage}}\end{center}

\bigskip

\begin{abstract}
Let $\A_n$ be Postnikov's affine nil-Temperley--Lieb algebra and let
$\mathcal R_e(t)=\sum_{J,\varepsilon}t^{|\varepsilon|}a_{J,\varepsilon}$ be the
height-graded cylindric $e$-ribbon adder, realised inside $\A_n[t]$ in
\cite{Clio201}.  We prove the cylindric splitting identity
\[
\sum_{b=0}^{e}t^{b}\,\bh_{e-b}\be_{b}
\;=\;\sum_{|I|=e}\ \sum_{\varepsilon}t^{|\varepsilon|}(1+t)^{r(I)}a_{I,\varepsilon}
\qquad(1\le e\le n-1),
\]
where $r(I)$ is the number of cyclic runs of $I\subseteq\ZZ/n\ZZ$.  The proof is
two observations about heaps: a horizontal-strip monomial \emph{empties} its own
support while a vertical-strip monomial \emph{requires} its support to be
occupied, so the two supports must be disjoint; and once they are disjoint, the
commutation class of the product is determined by which part each
\emph{edge-indexing} site lies in --- leaving the last site of each run free.
That free binary choice is the factor $(1+t)$.  Differentiating at $t=-1$ kills
every $I$ with two or more runs and yields the closed formula
$\mathcal R_e(-1)=\sum_{b\ge1}(-1)^{b-1}b\,\bh_{e-b}\be_b$, which exhibits
$\mathcal R_e(-1)$ as a quantum multiplication operator and proves the cylindric
Newton identity $r\,\bh_r=\sum_{e\le r}\mathcal R_e(-1)\bh_{r-e}$ for
$r\le n-1$, together with its defect $(-1)^{k-1}(n-k)q$ at $r=n$.  We also settle
the boundary cases $k\in\{1,n-1\}$: there $\mathcal R_e(t)$ is a \emph{monomial}
in $t$, so no value of $t$ is distinguished and the rigidity question is vacuous.
\end{abstract}

\section{Conventions, imported verbatim from \cite{Clio201}}\label{sec:conv}

Throughout, $n\ge2$ and $\A_n$ is the associative $\ZZ$-algebra on generators
$a_i$, $i\in\ZZ/n\ZZ$, with $a_ia_i=a_ia_{i+1}a_i=a_{i+1}a_ia_{i+1}=0$ and
$a_ia_j=a_ja_i$ for $i-j\not\equiv\pm1$.  It acts on
$\bigoplus_k\QH(\Gr_{kn})$ by $a_i=E_{i,i+1}$ for $0\le i\le n-2$ and
$a_{n-1}=q\,E_{n-1,0}$, where $E_{ij}$ moves a bead from site $i$ to site $j$
when $i$ is occupied and $j$ is empty and kills the basis vector otherwise.  A
partition $\lambda\in P_{kn}$ is the bead set
$S=\{\lambda_j+k-j:1\le j\le k\}\subseteq\ZZ/n\ZZ$, $|S|=k$.  Words are read in
\emph{application} order: the leftmost factor acts first
(\cite[\texttt{rem:conv}]{Clio201}; this is gap (G4) there and it is load-bearing here
too, since every statement below is about application order).

Call $I\subseteq\ZZ/n\ZZ$ \emph{proper} if $I\ne\ZZ/n\ZZ$.  For proper $I$, the
subgraph of the $n$-cycle induced on $I$ is a disjoint union of paths; its
maximal paths are the \emph{runs} of $I$, of which there are $r(I)$, and
$\mathrm{Ed}(I)$ is its edge set, $|\mathrm{Ed}(I)|=|I|-r(I)$.  We index each
edge by its smaller endpoint \emph{inside its run}, so that
\begin{equation}\label{eq:edgeindex}
\mathrm{Ed}(I)\ \longleftrightarrow\ I\setminus L(I),
\qquad
L(I):=\{\text{last site of each run of }I\},\quad |L(I)|=r(I),
\end{equation}
the edge indexed by $s$ being $\{s,s+1\}$.  For
$\varepsilon\in\{0,1\}^{\mathrm{Ed}(I)}$, $a_{I,\varepsilon}\in\A_n$ is the
product of $\{a_i:i\in I\}$, each used once, in any linear order for which
$\varepsilon_s=0$ means $a_s$ is applied before $a_{s+1}$ and $\varepsilon_s=1$
means $a_{s+1}$ is applied before $a_s$ (\cite[\texttt{def:aIeps}]{Clio201}).  This is
well defined, and $\varepsilon\mapsto a_{I,\varepsilon}$ is a bijection onto the
monomials of $\A_n$ with support $I$ using each generator once.  Postnikov's
elements are
\begin{equation}\label{eq:heDef}
\bh_m=\sum_{|I|=m,\ I\text{ proper}}a_{I,\mathbf 0},
\qquad
\be_b=\sum_{|I|=b,\ I\text{ proper}}a_{I,\mathbf 1},
\qquad \bh_0=\be_0=1,
\end{equation}
and $\mathcal R_e(t)=\sum_{|J|=e,\ r(J)=1}\sum_\varepsilon
t^{|\varepsilon|}a_{J,\varepsilon}$ is the sum over \emph{intervals}
(\cite[\texttt{eq:Rdef}]{Clio201}).  Finally we recall
\cite[\texttt{lem:run}]{Clio201}: on a run $J=[i,i+\ell-1]$ with orientation
$\varepsilon$, the monomial $a_{J,\varepsilon}$ is nonzero on the bead state $S$
if and only if
\begin{equation}\label{eq:pattern}
i\in S,\qquad i+\ell\notin S,\qquad
\text{and }\ i+s\in S\iff\varepsilon_s=1\quad(1\le s\le \ell-1),
\end{equation}
in which case it moves the bead at $i$ to $i+\ell$ and collects
$q^{[\,n-1\in J\,]}$.

\begin{lemma}[runs act independently]\label{lem:runsdisjoint}
Let $I$ be proper with runs $J_1,\dots,J_{r}$, $J_a=[i_a,i_a+\ell_a-1]$.  The
\emph{closed} intervals $\bar J_a:=[i_a,i_a+\ell_a]$ are pairwise disjoint.
Consequently $a_{I,\varepsilon}=\prod_a a_{J_a,\varepsilon|_{J_a}}$ in any
order, and it is nonzero on $S$ iff \eqref{eq:pattern} holds for every run
simultaneously.
\end{lemma}

\begin{proof}
By maximality $i_a+\ell_a\notin I$, so the next site of $I$ after
$i_a+\ell_a-1$ is at least $i_a+\ell_a+1$; hence the next run starts at
$\ge i_a+\ell_a+1$ and $\bar J_a\cap\bar J_{a'}=\varnothing$.  Generators
indexed in different runs are non-adjacent, hence commute; and
$a_s$ for $s\in J_a$ only reads and writes sites in $\bar J_a$.
\end{proof}

\section{The two lemmas}\label{sec:two}

Everything rests on the following pair.  Neither is about Young diagrams: they
are statements about occupancy and about the trace monoid.

\begin{lemma}[a horizontal monomial empties its support; a vertical monomial
fills its own]\label{lem:emptyfill}
Let $I$ be proper and $S$ a bead state.
\begin{enumerate}
\item[(a)] If $a_{I,\mathbf 0}\,\sigma_S\ne0$ and $U$ is the resulting bead
state, then $I\cap U=\varnothing$.
\item[(b)] If $a_{I,\mathbf 1}\,\sigma_S\ne0$ then $I\subseteq S$.
\end{enumerate}
\end{lemma}

\begin{proof}
Both are read off \eqref{eq:pattern} run by run, legitimately by
Lemma~\ref{lem:runsdisjoint}.

(a) For a run $J=[i,i+\ell-1]$ with $\varepsilon=\mathbf 0$, \eqref{eq:pattern}
says $i\in S$ and $i+s\notin S$ for $1\le s\le\ell$.  The move deletes the bead
at $i$ and creates one at $i+\ell$, so in $U$ the sites $i,i+1,\dots,i+\ell-1$
are all empty --- that is, $J\cap U=\varnothing$.  Other runs do not touch
$\bar J$.  Taking the union over runs gives $I\cap U=\varnothing$.

(b) For a run $J=[i,i+\ell-1]$ with $\varepsilon=\mathbf 1$, \eqref{eq:pattern}
says $i\in S$ and $i+s\in S$ for $1\le s\le\ell-1$, i.e.\ $J\subseteq S$.
\end{proof}

\begin{corollary}[no repeated generator survives]\label{cor:norepeat}
If $I'\cap I''\ne\varnothing$ then $a_{I',\mathbf 0}\,a_{I'',\mathbf 1}=0$ as an
operator on $\bigoplus_k\QH(\Gr_{kn})$.
\end{corollary}

\begin{proof}
Suppose $a_{I',\mathbf 0}a_{I'',\mathbf 1}\sigma_S\ne0$ for some $S$ (recall
$a_{I',\mathbf 0}$ acts first).  Let $U$ be the state after the first factor.
By Lemma~\ref{lem:emptyfill}(a), $I'\cap U=\varnothing$; by
Lemma~\ref{lem:emptyfill}(b) applied at $U$, $I''\subseteq U$.  Hence
$I'\cap I''\subseteq I'\cap U=\varnothing$.
\end{proof}

\begin{lemma}[a disjoint split names an orientation, and only the run-ends are
free]\label{lem:split}
Let $I$ be proper and let $I=I'\sqcup I''$.  Then
\[
a_{I',\mathbf 0}\;a_{I'',\mathbf 1}\;=\;a_{I,\varepsilon(I'')}
\qquad\text{in }\A_n,
\qquad\text{where }\ \varepsilon(I'')_s=\begin{cases}1,&s\in I'',\\0,&s\in I',\end{cases}
\]
for every edge index $s\in\mathrm{Ed}(I)$.  Conversely, for each
$\varepsilon\in\{0,1\}^{\mathrm{Ed}(I)}$ the splits $I=I'\sqcup I''$ with
$\varepsilon(I'')=\varepsilon$ are exactly those obtained by putting
$s\in I''$ for every edge index $s$ with $\varepsilon_s=1$, $s\in I'$ for every
edge index $s$ with $\varepsilon_s=0$, and assigning each of the $r(I)$ sites of
$L(I)$ freely to $I'$ or to $I''$.  Hence
\begin{equation}\label{eq:twochoices}
\sum_{\substack{I=I'\sqcup I''}}t^{|I''|}\,a_{I',\mathbf 0}a_{I'',\mathbf 1}
\;=\;\sum_{\varepsilon}t^{|\varepsilon|}(1+t)^{r(I)}\,a_{I,\varepsilon}.
\end{equation}
\end{lemma}

\begin{proof}
The concatenated word $w$ uses each $a_s$, $s\in I$, exactly once, so it is a
representative of $a_{I,\varepsilon}$ for exactly one
$\varepsilon\in\{0,1\}^{\mathrm{Ed}(I)}$ --- provided we check that any two
linear orders agreeing on all edges of $I$ give the same element of $\A_n$.
They do: two linear extensions of the partial order generated by the edge
constraints differ by transpositions of adjacent-in-word incomparable elements,
and two incomparable elements of $I$ are non-adjacent in $\ZZ/n\ZZ$ (if $u,u+1$
both lie in $I$ then $\{u,u+1\}$ is an edge of $I$ and they are comparable),
hence commute.

It remains to identify $\varepsilon$.  Fix an edge $\{s,s+1\}$ of $I$ and
consider the four cases.  If $s\in I'$ and $s+1\in I'$, then $s,s+1$ lie in the
same run of $I'$ (they are adjacent and both in $I'$), and $\mathbf 0$ puts
$a_s$ before $a_{s+1}$: $\varepsilon_s=0$.  If $s\in I'$ and $s+1\in I''$, then
$a_s$ is in the first factor and $a_{s+1}$ in the second, so $a_s$ is applied
first: $\varepsilon_s=0$.  If $s\in I''$ and $s+1\in I'$, then $a_{s+1}$ is in
the first factor and $a_s$ in the second, so $a_{s+1}$ is applied first:
$\varepsilon_s=1$.  If $s\in I''$ and $s+1\in I''$, they lie in the same run of
$I''$ and $\mathbf 1$ puts $a_{s+1}$ before $a_s$: $\varepsilon_s=1$.  In all
four cases $\varepsilon_s=[\,s\in I''\,]$, which is the assertion.

The converse is immediate: the map $I''\mapsto\varepsilon(I'')$ reads off the
membership of the $|I|-r(I)$ edge-indexing sites and forgets the membership of
the $r(I)$ sites of $L(I)$, which by \eqref{eq:edgeindex} is the whole of the
remaining data.  Each fibre therefore has $2^{r(I)}$ elements, one for each
subset $L''\subseteq L(I)$, and $|I''|=|\varepsilon|+|L''|$.  Summing
$t^{|I''|}$ over the fibre gives $t^{|\varepsilon|}(1+t)^{r(I)}$, which is
\eqref{eq:twochoices}.
\end{proof}

\begin{remark}[what happened to the classical proof]\label{rem:classical}
Lemma~\ref{lem:split} is the cylindric replacement for three lemmas of
\cite{Clio75}: ``a $2\times2$ kills every splitting'', ``splittings factor over
components'', and ``a ribbon has exactly two splittings''.  All three are
absorbed.  The $2\times2$ obstruction is invisible here because a bead move of
the form \eqref{eq:pattern} can never contain one --- the monomials of $\A_n$
with a given support \emph{are} the $2\times2$-free shapes, by
\cite[\texttt{thm:A}]{Clio201}; the factorisation over components is
Lemma~\ref{lem:runsdisjoint}; and ``exactly two splittings'' is the free choice
of the single site $L(J)$ of a one-run $J$.  In the classical argument the
factor $(1+t)$ emerges after a page of interlacing inequalities
(\cite[\texttt{lem:ribbon}]{Clio75}) as $\alpha_1\in\{a_1-1,a_1\}$; here it is the last
generator of the run, which may be applied in either factor because no
constraint mentions it.  The two are the same fact: $\alpha_1$ is the number of
cells of the first row placed in the horizontal half, and the run-end site is
the corresponding generator.
\end{remark}

\section{The splitting identity and its derivative}

\begin{theorem}[cylindric splitting identity]\label{thm:C1}
For $1\le e\le n-1$, as operators on $\bigoplus_k\QH(\Gr_{kn})$,
\begin{equation}\label{eq:C1}
\sum_{b=0}^{e}t^{b}\,\bh_{e-b}\,\be_{b}
\;=\;
\sum_{\substack{I\subseteq\ZZ/n\ZZ\\ |I|=e}}\ \sum_{\varepsilon\in\{0,1\}^{\mathrm{Ed}(I)}}
t^{|\varepsilon|}\,(1+t)^{r(I)}\,a_{I,\varepsilon}.
\end{equation}
\end{theorem}

\begin{proof}
Since $e\le n-1$, every $I$ with $|I|=e$ is proper, and so is every subset of
it; thus all the monomials occurring below are among those of
\eqref{eq:heDef}.  Expanding by \eqref{eq:heDef},
\[
\sum_{b=0}^{e}t^{b}\bh_{e-b}\be_{b}
=\sum_{b=0}^{e}t^{b}\sum_{|I'|=e-b}\sum_{|I''|=b}a_{I',\mathbf 0}a_{I'',\mathbf 1}.
\]
By Corollary~\ref{cor:norepeat} every term with $I'\cap I''\ne\varnothing$
vanishes, so the double sum is over ordered pairs of \emph{disjoint} sets with
$|I'\sqcup I''|=e$.  Grouping these by $I:=I'\sqcup I''$ and applying
\eqref{eq:twochoices} to each $I$ gives \eqref{eq:C1}.
\end{proof}

\begin{corollary}[$t=-1$ recovers Postnikov's vanishing]\label{cor:posvanish}
For $1\le e\le n-1$, $\sum_{b=0}^{e}(-1)^b\bh_{e-b}\be_b=0$; that is, the
coefficients of $t^1,\dots,t^{n-1}$ in $\bigl(1+\sum_i\be_it^i\bigr)
\bigl(1+\sum_j\bh_j(-t)^j\bigr)$ vanish.
\end{corollary}

\begin{proof}
$r(I)\ge1$ for $|I|=e\ge1$, so every term on the right of \eqref{eq:C1} carries
a factor $(1+t)$.
\end{proof}

This is the part of Postnikov's Lemma (\texttt{lem:e-h-commute-in-AN},
\cite[ll.2882--2893]{Pos}) that lives below degree $n$; it is recovered here,
not assumed.  The new content is one derivative further.

\begin{theorem}[closed formula for the ribbon adder at $t=-1$]\label{thm:pe}
For $1\le e\le n-1$, as operators on $\bigoplus_k\QH(\Gr_{kn})$,
\begin{equation}\label{eq:pe}
\boxed{\;\mathcal R_e(-1)\;=\;\sum_{b=1}^{e}(-1)^{b-1}\,b\;\bh_{e-b}\,\be_{b}.\;}
\end{equation}
\end{theorem}

\begin{proof}
Both sides of \eqref{eq:C1} are polynomials in $t$ with operator coefficients;
differentiate in $t$ and set $t=-1$.

On the right, the $I$-th summand is
$\sum_\varepsilon t^{|\varepsilon|}(1+t)^{r(I)}a_{I,\varepsilon}$, whose
derivative is
$\sum_\varepsilon\bigl(|\varepsilon|t^{|\varepsilon|-1}(1+t)^{r(I)}
+r(I)t^{|\varepsilon|}(1+t)^{r(I)-1}\bigr)a_{I,\varepsilon}$.  At $t=-1$ the
first term vanishes because $r(I)\ge1$, and the second vanishes unless
$r(I)=1$, in which case it equals
$\sum_\varepsilon(-1)^{|\varepsilon|}a_{I,\varepsilon}$.  Now $|I|=e$ and
$r(I)=1$ says exactly that $I$ is an interval of length $e$, so the surviving
sum is $\sum_{|J|=e,\,r(J)=1}\sum_\varepsilon(-1)^{|\varepsilon|}a_{J,\varepsilon}
=\mathcal R_e(-1)$ by \cite[\texttt{thm:A}]{Clio201}.

On the left, $\frac{d}{dt}\sum_b t^b\bh_{e-b}\be_b
=\sum_{b\ge1}b\,t^{b-1}\bh_{e-b}\be_b$, which at $t=-1$ is
$\sum_{b\ge1}(-1)^{b-1}b\,\bh_{e-b}\be_b$.
\end{proof}

\begin{remark}
The same derivative, performed in $\Lam$ rather than in $\A_n$, is the classical
identity $p_e=\sum_{b=1}^{e}(-1)^{b-1}b\,h_{e-b}e_b$: from $H(z)E(-z)=1$ one
gets $\frac{\partial}{\partial t}H(z)E(tz)\bigr|_{t=-1}=zE'(-z)/E(-z)
=z\,\frac{d}{dz}\log H(z)=\sum_{e\ge1}p_ez^{e}$.  So \eqref{eq:pe} is literally
the statement that the classical expression of $p_e$ in the $h$'s and $e$'s is
valid verbatim on the cylinder in degrees $\le n-1$ --- which is why no wrap
term can appear there.  The wrap is confined to $e=n$, where $I=\ZZ/n\ZZ$ is not
proper and Theorem~\ref{thm:C1} has nothing to say.
\end{remark}

\section{The cylindric Newton identity}

\begin{lemma}[the strip adders are the quantum multiplications, in the full
range]\label{lem:mult}
Fix $1\le k\le n-1$ and write $M_x$ for quantum multiplication by
$x\in\QH(\Gr_{kn})$.  Then on $\QH(\Gr_{kn})$,
\[
\bh_j=M_{\sigma_j}\quad(0\le j\le n-1),
\qquad
\be_b=M_{\sigma_{1^b}}\quad(0\le b\le n-1),
\]
with the convention $\sigma_j=0$ for $j>n-k$ and $\sigma_{1^b}=0$ for $b>k$.
\end{lemma}

\begin{proof}
For $1\le j\le n-k$ and $1\le b\le k$ this is Postnikov's quantum Pieri
corollary \cite[\texttt{cor:pieri-AN}, ll.2935--2947]{Pos}, read at source; the
cases $j=0,b=0$ are trivial.  It remains to see that the operators vanish
outside those ranges, which the bead model makes immediate.

If $b>k$: by Lemma~\ref{lem:emptyfill}(b) a nonzero $a_{I,\mathbf 1}$ with
$|I|=b$ requires $I\subseteq S$, impossible since $|S|=k<b$.  Hence
$\be_b=0=M_{\sigma_{1^b}}$.

If $j>n-k$: let $I$ be proper with $|I|=j$ and $a_{I,\mathbf 0}\sigma_S\ne0$.
For each run $J_a=[i_a,i_a+\ell_a-1]$, \eqref{eq:pattern} with
$\varepsilon=\mathbf 0$ forces the $\ell_a$ sites
$i_a+1,\dots,i_a+\ell_a$ to be empty in $S$.  These sets are pairwise disjoint
over the runs, being contained in the pairwise disjoint $\bar J_a$ of
Lemma~\ref{lem:runsdisjoint} and each of size $\ell_a$; so $S$ has at least
$\sum_a\ell_a=j$ empty sites.  But $S$ has exactly $n-k<j$ empty sites.  Hence
$\bh_j=0=M_{\sigma_j}$.
\end{proof}

\begin{theorem}[Newton, and the defect]\label{thm:newton}
Let $1\le k\le n-1$ and put $p_e:=\mathcal R_e(-1)$ for $1\le e\le n-1$.  Then,
as operators on $\QH(\Gr_{kn})$:
\begin{enumerate}
\item[(i)] each $p_e$ is a quantum multiplication operator, namely
  $p_e=M_{\varphi(p_e)}$, where $\varphi:\Lam[q]\to\QH(\Gr_{kn})$ is the ring
  homomorphism with $\varphi(e_i)=\sigma_{1^i}$; in particular the $p_e$ commute
  with one another and with every $\bh_j$;
\item[(ii)] $\displaystyle r\,\bh_r=\sum_{e=1}^{r}p_e\,\bh_{r-e}$ for
  $1\le r\le n-1$;
\item[(iii)] $\displaystyle\sum_{e=1}^{n-1}p_e\,\bh_{n-e}
  =(-1)^{k-1}(n-k)\,q\cdot\mathrm{Id}$.
\end{enumerate}
\end{theorem}

\begin{proof}
$\varphi$ is well defined because $\Lam=\ZZ[e_1,e_2,\dots]$ is a free polynomial
ring, and we extend it $\ZZ[q]$-linearly.  We first record
\begin{equation}\label{eq:phih}
\varphi(h_j)=\sigma_j\qquad(0\le j\le n-1).
\end{equation}
Indeed $\sum_{b=0}^{j}(-1)^bh_{j-b}e_b=0$ in $\Lam$ for $j\ge1$; applying
$\varphi$ gives $\sum_{b=0}^{j}(-1)^b\varphi(h_{j-b})\sigma_{1^b}=0$.  On the
other side, the coefficient of $t^j$ in Postnikov's presentation
$e(t)h(-t)=1+(-1)^{n-k}qt^n$ \cite[\texttt{eq:QH}, ll.500--507]{Pos} is
$\sum_{b=0}^{j}(-1)^{j-b}\sigma_{1^b}\sigma_{j-b}=0$ for $1\le j\le n-1$, i.e.\
the same recursion.  Both sequences start at $\varphi(h_0)=1=\sigma_0$, and the
recursion determines the $j$-th term from the earlier ones (the $b=0$ term is
$\varphi(h_j)$ resp.\ $\sigma_j$, with coefficient $\pm1$), so
\eqref{eq:phih} follows by induction.

(i)  By Theorem~\ref{thm:pe} and Lemma~\ref{lem:mult},
\[
p_e=\sum_{b=1}^{e}(-1)^{b-1}b\,\bh_{e-b}\be_b
=\sum_{b=1}^{e}(-1)^{b-1}b\,M_{\sigma_{e-b}}M_{\sigma_{1^b}}
=M_{\sum_{b=1}^{e}(-1)^{b-1}b\,\sigma_{e-b}*\sigma_{1^b}},
\]
using that $\QH(\Gr_{kn})$ is a commutative associative ring, so
$M_xM_y=M_{x*y}$.  By \eqref{eq:phih} and the classical identity
$p_e=\sum_{b\ge1}(-1)^{b-1}b\,h_{e-b}e_b$ in $\Lam$ (Remark after
Theorem~\ref{thm:pe}), the subscript is $\varphi(p_e)$.  Multiplication
operators of a commutative ring commute with one another, and $\bh_j=M_{\sigma_j}$.

(ii)  For $1\le r\le n-1$, by (i), Lemma~\ref{lem:mult} and \eqref{eq:phih},
\[
\sum_{e=1}^{r}p_e\bh_{r-e}
=\sum_{e=1}^{r}M_{\varphi(p_e)}M_{\varphi(h_{r-e})}
=M_{\varphi\left(\sum_{e=1}^{r}p_eh_{r-e}\right)}
=M_{\varphi(r\,h_r)}
=r\,M_{\sigma_r}=r\,\bh_r,
\]
the third equality being Newton's identity in $\Lam$.

(iii)  The same computation at $r=n$, where $\bh_{n-e}=M_{\varphi(h_{n-e})}$
still holds for $1\le e\le n-1$ because then $n-e\le n-1$, gives
\[
\sum_{e=1}^{n-1}p_e\bh_{n-e}
=M_{\varphi\left(\sum_{e=1}^{n-1}p_eh_{n-e}\right)}
=M_{\varphi(n\,h_n-p_n)}
=M_{n\varphi(h_n)-\varphi(p_n)} .
\]
It remains to evaluate the two scalars.  Taking $j=n$ in the $\Lam$-recursion,
$\varphi(h_n)=-\sum_{b=1}^{n}(-1)^b\sigma_{1^b}\varphi(h_{n-b})
=-\sum_{b=1}^{n}(-1)^b\sigma_{1^b}\sigma_{n-b}$ by \eqref{eq:phih}; and the
$t^n$ coefficient of $e(t)h(-t)=1+(-1)^{n-k}qt^n$ reads
$\sum_{b=0}^{n}(-1)^{n-b}\sigma_{1^b}\sigma_{n-b}=(-1)^{n-k}q$, whose $b=0$ term
is $\sigma_n=0$ (as $n>n-k$); multiplying by $(-1)^n$ gives
$\sum_{b=1}^{n}(-1)^{b}\sigma_{1^b}\sigma_{n-b}=(-1)^{k}q$.  Hence
$\varphi(h_n)=(-1)^{k-1}q$.

For $\varphi(p_n)=\sum_{b=1}^{n}(-1)^{b-1}b\,\varphi(h_{n-b})\sigma_{1^b}$, the
factor $\sigma_{1^b}$ vanishes for $b>k$ and, for $1\le b\le n-1$,
$\varphi(h_{n-b})=\sigma_{n-b}$ vanishes unless $n-b\le n-k$, i.e.\ $b\ge k$;
the term $b=n$ carries $\sigma_{1^n}=0$ since $n>k$.  So only $b=k$ survives:
$\varphi(p_n)=(-1)^{k-1}k\,\sigma_{n-k}*\sigma_{1^k}=(-1)^{k-1}k\,q$, using the
single quantum relation $e_kh_{n-k}=q$ \cite[\texttt{thm:C}(ii)]{Clio201}.  Therefore
\[
n\varphi(h_n)-\varphi(p_n)=n(-1)^{k-1}q-(-1)^{k-1}kq=(-1)^{k-1}(n-k)q. \qedhere
\]
\end{proof}

\begin{corollary}[quantum Murnaghan--Nakayama, unconditionally]\label{cor:MN}
For $1\le e\le n-1$ and $\lambda\in P_{kn}$,
\[
\sum_{\substack{\mu:\ \mu/\lambda\ \text{a cylindric }e\text{-ribbon}}}
(-1)^{\hgt(\mu/\lambda)}q^{\,w(\mu/\lambda)}\,\sigma_\mu
\;=\;\varphi(p_e)*\sigma_\lambda ,
\]
where $w$ counts seam crossings.  In particular
$[\mathcal R_e(-1),\bh_j]=0$ for all $e,j$.
\end{corollary}

\begin{proof}
The left side is $\mathcal R_e(-1)\sigma_\lambda$ by
\cite[\texttt{eq:ribbondef} and \texttt{thm:A}]{Clio201}; apply Theorem~\ref{thm:newton}(i).
\end{proof}

\noindent
This is \cite[\texttt{prop:MN}]{Clio201}, whose proof there assumed
Theorem \texttt{thm:B}(ii); the dependence is now reversed and the assumption
discharged.

\section{The boundary cases $k=1$ and $k=n-1$ (gap (G3))}

\begin{proposition}\label{prop:G3}
Let $1\le e\le n-1$.
\begin{enumerate}
\item[(a)] On $\QH(\Gr_{1,n})$, $\mathcal R_e(t)=\mathcal R_e(0)$: the operator
  does not depend on $t$ at all.
\item[(b)] On $\QH(\Gr_{n-1,n})$, $\mathcal R_e(t)=t^{\,e-1}\mathcal R_e(1)$:
  the operator is $t^{e-1}$ times a $t$-free operator.
\end{enumerate}
Consequently, at $k\in\{1,n-1\}$ the identity $[\mathcal R_e(t),\bh_j]=0$ holds
either for every $t$ or for no $t$, and it does hold for every $t$ by
Corollary~\ref{cor:MN}.  The uniqueness assertion of
\cite[\texttt{thm:B}(i)]{Clio201} is therefore not merely untested at the boundary: it
is \emph{vacuous} there, because the $t$-line has collapsed to a single
monomial.  This is a rank degeneracy of the Hall--Littlewood grading, not a
statement about $\mathbb P^{n-1}$.
\end{proposition}

\begin{proof}
Recall from \eqref{eq:pattern} (equivalently \cite[\texttt{eq:ribbondef}]{Clio201}) that the
$(S,j)$ term of $\mathcal R_e(t)$ carries $t^{\#\{1\le s\le e-1:\ j+s\in S\}}$.

(a) $|S|=1$, say $S=\{\beta\}$, and the only legal source is $j=\beta$.  For
$1\le s\le e-1\le n-2$ we have $\beta+s\not\equiv\beta$, so $\beta+s\notin S$
and the exponent is $0$.

(b) $|S|=n-1$; let $g$ be the unique empty site.  A move $j\mapsto j+e$ is legal
iff $j+e\notin S$, i.e.\ $j+e=g$.  For $1\le s\le e-1$, $j+s\equiv g=j+e$ would
force $s\equiv e\pmod n$, impossible; so $j+s\ne g$, i.e.\ $j+s\in S$, and the
exponent is $e-1$ for every legal term.
\end{proof}

\begin{remark}
(a) and (b) are exchanged by the duality $\rho$ of
\cite[\texttt{cor:duality}]{Clio201}, $\rho\mathcal R_e(t)\rho^{-1}=t^{e-1}\mathcal R_e(1/t)$,
which is consistent: a $t$-free operator is carried to a $t^{e-1}$-multiple of a
$t$-free one.
\end{remark}

\section{Computational verification}\label{sec:verif}

All checks are exact and symbolic (SymPy over $\ZZ[q,t]$), in
\texttt{proofs/code-2026-09-24/}, built on the $\A_n$ model of
\texttt{proofs/code-q237/anTL.py}, which implements $a_i$ \emph{as the bead
move}, so that the nil-Temperley--Lieb relations are consequences of the model rather
than inputs to it.  Each check below tests a \emph{step} of the proof, not its
conclusion; the conclusion (Theorem~\ref{thm:newton}(ii)--(iii)) was already
verified for $3\le n\le7$ in \cite{Clio201} and was deliberately not re-run.

\begin{enumerate}
\item \textbf{Corollary~\ref{cor:norepeat}} (\texttt{check1\_no\_repeat.py}).
For $2\le n\le6$, all $0\le k\le n$, and all ordered pairs of nonempty proper
$I',I''$ with $I'\cap I''\ne\varnothing$: $a_{I',\mathbf 0}a_{I'',\mathbf 1}=0$.
$27\,846$ pairs, \textbf{0 violations}.
\item \textbf{Lemma~\ref{lem:split}} (\texttt{check4\_split\_orientation.py}).
For $2\le n\le7$, all $k$, all proper $I$ and all splits $I=I'\sqcup I''$:
$a_{I',\mathbf 0}a_{I'',\mathbf 1}=a_{I,\varepsilon}$ with
$\varepsilon_s=[s\in I'']$.  $22\,776$ splits, \textbf{0 violations}.
\item \textbf{Theorem~\ref{thm:C1}} (\texttt{check2\_C1\_identity.py}).
For $2\le n\le7$, all $0\le k\le n$, all $1\le e\le n-1$: the two sides of
\eqref{eq:C1} agree entrywise as polynomials in $\ZZ[q,t]$.
\textbf{0 violations}.
\item \textbf{Theorem~\ref{thm:pe}} (\texttt{check3\_pe\_formula.py}).
Same ranges: $\sum_{b\ge1}(-1)^{b-1}b\,\bh_{e-b}\be_b$ equals the independent
bead definition of $\mathcal R_e(t)$ specialised at $t=-1$.
\textbf{0 violations}.
\item \textbf{Lemma~\ref{lem:mult}, vanishing half}
(\texttt{check5\_vanishing.py}).  For $2\le n\le8$, all $k$:
$\be_b=0$ for $b>k$ and $\bh_j=0$ for $j>n-k$.  \textbf{0 violations}.
\item \textbf{Proposition~\ref{prop:G3}} (\texttt{g3\_boundary.py}).
For $3\le n\le8$ and $k\in\{1,n-1\}$, all $1\le e\le n-1$: the set of
$t$-degrees occurring in the matrix of $\mathcal R_e(t)$ is $\{0\}$ when $k=1$
and $\{e-1\}$ when $k=n-1$.
\end{enumerate}

\section{Gaps}\label{sec:gaps}

\begin{enumerate}
\item[(H1)] \textbf{Uniqueness of $t=-1$ is still open.}  We prove
$[\mathcal R_e(-1),\bh_j]=0$ (Corollary~\ref{cor:MN}); we do \emph{not} prove
that $t=-1$ is the only such point for $2\le k\le n-2$, which is the remaining
half of \cite[\texttt{thm:B}(i)]{Clio201} and is verified only for $4\le n\le6$.
Theorem~\ref{thm:C1} does not settle it: it shows that
$\sum_{c\ge1}(1+t)^cA_c(t)$ is a multiplication operator for \emph{every} $t$,
where $A_c(t)$ collects the $I$ with $r(I)=c$, and one cannot isolate
$A_1=\mathcal R_e$ from that.  Proposition~\ref{prop:G3} shows the question is
vacuous at $k\in\{1,n-1\}$, which is why the hypothesis $2\le k\le n-2$ is the
right one, but it does not prove uniqueness inside that range.
\item[(H2)] \textbf{The route predicted by the brief was not taken, and its
prediction is untested.}  The WAKE brief for this session conjectured that the
defect at $r=n$ would be exhibited by enumerating wrap-around contributions
indexed by the $n-k$ \emph{gaps} of the bead configuration, and proposed that as
a falsifiable test.  The proof above never enumerates wrap contributions: the
$n-k$ arises as $n\cdot1-k\cdot1$, from $n\varphi(h_n)$ against the single
surviving index $b=k$ in $\varphi(p_n)$.  The gap-indexed picture may still be
true; \emph{it is neither used nor checked here}, and nothing in this note is
evidence for it.
\item[(H3)] \textbf{Lemma~\ref{lem:mult} rests on a reading of Postnikov.}
\texttt{cor:pieri-AN} is stated at \cite[ll.2935--2947]{Pos} for $i=1,\dots,k$
and $j=1,\dots,n-k$ only; the extension to the full range $0\le\cdot\le n-1$ is
the vanishing argument given here, which is ours.  The line numbers are internal
locators in the e-print I hold, not published theorem numbers.  The
application-order convention of \cite[\texttt{rem:conv}]{Clio201} (gap (G4) there)
remains a reading of ll.2839--2877, not a quotation, and Lemma~\ref{lem:split}
is entirely a statement about application order --- if that reading is wrong,
the four-case analysis in its proof inverts and $\varepsilon_s=[s\in I']$
instead.  The identity \eqref{eq:C1} would then hold with $\bh$ and $\be$
exchanged, which is the $t\mapsto1/t$ image of the same statement; nothing
downstream changes.
\item[(H4)] \textbf{Two of the three novelty findings are agent-reported, not
read by me.}  I verified Korff \cite{Korff19} \texttt{lem:cylMNrule} at source
today.  I have \emph{not} opened Morrison--Sottile \cite{MS} (index entry:
\texttt{agent-summary}) or Korff \cite{Korff18} \texttt{lem:CR}.  A sub-agent
summary is a hypothesis; these two are recorded as reported, and the paper
claims nothing that depends on their being right --- it would only become
\emph{more} novel if they were wrong, which is the safe direction.  What is
genuinely owed is the convention reconciliation with \cite{Korff19}: my
$(-1)^{k-1}(n-k)q$ against his $(-1)^k(n-k)$, under conjugation to
$\mathcal P^+_{n-k,n}$ and normalisation by $(t-1)^{\ell(\nu)}$.  Until that is
done I do not assert that the two constants are the same constant.
\item[(H4$'$)] \textbf{Process failure, recorded because it nearly cost the
session.}  The BROWSE gate sat at the head of \texttt{memory/SUMMARY.md} with
the words ``READ BEFORE PROVING''.  I read the PROVE brief, which pointed at it,
and went straight to the mathematics; I found the gate only when updating
\texttt{SUMMARY.md} at the end.  Had the proof gone the other way --- had I
spent the session on the \emph{conclusion} rather than the \emph{route} --- the
whole of it would have been wasted, and the brief said exactly that
(``Instrument the ROUTE'').  The outcome is acceptable only by luck.
\item[(H5)] \textbf{Gap (G2) of \cite{Clio201} is moot, not resolved.}  That
gap proposed proving the $q=0$ case first and lifting.  The present proof is
uniform in $q$ and never performs that reduction, so the reduction remains
uninspected.  It is no longer on the critical path.
\end{enumerate}

\begin{thebibliography}{9}
\bibitem{Pos} A.~Postnikov, \emph{Affine approach to quantum Schubert calculus},
\texttt{arXiv:math/0205165v2}.  Cited by internal \LaTeX{} labels and line
numbers of the e-print source at
\texttt{sources/math0205165-postnikov/postnikov-affine-v2.tex}, read at source
on 2026-09-23 and 2026-09-24.
\bibitem{MS} A.~Morrison and F.~Sottile, \emph{Two Murnaghan--Nakayama rules in
Schubert calculus}, Ann.\ Comb.\ \textbf{22} (2018) 363--375,
\texttt{arXiv:1507.06569v2}.  \textbf{Not read by me};
reported by BROWSE 2026-09-24 c1 at \texttt{Th:two}, src ll.185--197.  Index
entry graded \texttt{agent-summary}.
\bibitem{Korff18} C.~Korff, \emph{Cylindric symmetric functions and positivity},
\texttt{arXiv:1804.05647}.  \texttt{lem:CR} (src ll.1925--1936) and
\texttt{cylSchur} (ll.1957--1960) reported by BROWSE; \textbf{those lines not
read by me}.
\bibitem{Korff19} C.~Korff, \emph{Cylindric Hecke characters and Gromov--Witten
invariants via the asymmetric six-vertex model}.  E-print source at
\texttt{sources/korff-cylindric-20260923/1906.02565-src/draft\_02.tex};
\texttt{lem:cylMNrule} read at source 2026-09-24, lines as cited above.
\bibitem{Clio201} Clio, \emph{The height of a ribbon is the orientation of a
heap: the Hall--Littlewood $t$-family inside the affine nil-Temperley--Lieb
algebra}, \texttt{proofs/2026-09-23-c2-ribbon-height-is-a-heap-orientation.tex},
commit \texttt{clio-vega/proofs@0a01bb9}.
\bibitem{Clio75} Clio, \emph{The multiplication defect of the Hall--Littlewood
ribbon family}, \texttt{proofs/2026-09-03-Q75-multiplication-defect.tex}.
\end{thebibliography}

\end{document}
